% =====================================================================
%  Thème 3 — Relation fondamentale & formules d'addition — NIVEAU DIFFICILE
%  Exercices
% =====================================================================
\documentclass[11pt,a4paper]{article}
\usepackage[utf8]{inputenc}
\usepackage[T1]{fontenc}
\usepackage{lmodern}
\usepackage[french]{babel}
\usepackage{amsmath,amssymb,mathtools}
\usepackage{icomma}
\usepackage{xcolor}
\usepackage{colortbl}
\usepackage{tikz}
\usepackage[most]{tcolorbox}
\usepackage{enumitem}
\usepackage{array}
\usepackage[a4paper,margin=2.2cm]{geometry}
\usetikzlibrary{arrows.meta,calc,angles,quotes}

\definecolor{primary}{RGB}{214,51,108}
\definecolor{primarydark}{RGB}{140,20,64}

\DeclareMathOperator{\tg}{tg}
\DeclareMathOperator{\cotg}{cotg}
\newcommand{\degr}{^{\circ}}
\newcommand{\Z}{\mathbb{Z}}

\newtcolorbox{consigne}{enhanced,breakable,colback=primary!5,colframe=primary,
  boxrule=0.8pt,arc=2pt,left=3mm,right=3mm,top=2mm,bottom=2mm}

\newcounter{exo}
\newcommand{\exo}[1]{\medskip\refstepcounter{exo}%
  \noindent{\large\bfseries\color{primarydark}Exercice \theexo}%
  \ \textcolor{primary}{\textbullet}\ \textbf{#1}\par\nopagebreak\smallskip}

\setlength{\parindent}{0pt}
\pagestyle{empty}

\begin{document}

\begin{center}
{\LARGE\bfseries\color{primary}Thème 3 — Relation fondamentale \& formules d'addition}\\[2pt]
{\color{primarydark}\textsc{Niveau Difficile — Exercices}}
\end{center}
\hrule height 1pt
\medskip

\begin{consigne}
\textbf{Objectif :} enchaîner \emph{relation fondamentale $\to$ duplication $\to$ addition} sur une même
donnée, et démontrer proprement une formule. Les sous-questions se construisent : chaque réponse sert à
la suivante. \textbf{Sans calculatrice}, valeurs \emph{exactes}.
\end{consigne}

\exo{La chaîne complète}
On suppose $\sin\alpha=\dfrac{1}{3}$ et $\alpha$ au \textbf{quadrant I}.
\begin{enumerate}[label=\textbf{(\alph*)},leftmargin=2em,itemsep=2pt]
\item Calcule $\cos\alpha$ (valeur exacte, signe justifié par le quadrant).
\item En déduire $\sin2\alpha$ et $\cos2\alpha$.
\item Calcule $\sin\!\Big(2\alpha-\dfrac{\pi}{3}\Big)$.
\item Calcule $\cos\!\Big(2\alpha+\dfrac{\pi}{3}\Big)$.
\end{enumerate}

\exo{Démontrer les formules de l'angle double}
On part uniquement de $\cos(a+b)=\cos a\cos b-\sin a\sin b$ et de $\cos^2a+\sin^2a=1$.
\begin{enumerate}[label=\textbf{(\alph*)},leftmargin=2em,itemsep=2pt]
\item En posant $b=a$, démontre que $\cos2a=\cos^2a-\sin^2a$.
\item En déduire les deux autres écritures : $\cos2a=2\cos^2a-1$ et $\cos2a=1-2\sin^2a$.
\item En déduire la formule de linéarisation $\sin^2x=\dfrac{1-\cos2x}{2}$.
\end{enumerate}

\exo{Point sur le cercle et angle double}
Sur le cercle trigonométrique, le point $M$ correspond à l'angle $\alpha$ situé au \textbf{quadrant II},
et l'on sait que $\tg\alpha=-\dfrac{3}{2}$.
\begin{center}
\begin{tikzpicture}[scale=1.7,line cap=round,>=Stealth]
  \draw[primary,line width=1pt] (0,0) circle (1);
  \draw[->,black!70] (-1.3,0)--(1.35,0) node[right]{\scriptsize$\cos$};
  \draw[->,black!70] (0,-1.3)--(0,1.35) node[above]{\scriptsize$\sin$};
  \coordinate (M) at (146.3:1);
  \draw[primarydark,thick] (0,0)--(M);
  \fill[primarydark] (M) circle (0.028) node[above left]{\scriptsize$M(\alpha)$};
  \draw[primarydark,->] (0.28,0) arc (0:146.3:0.28);
  \fill (0,0) circle (0.02) node[below right]{\scriptsize$O$};
\end{tikzpicture}
\end{center}
\begin{enumerate}[label=\textbf{(\alph*)},leftmargin=2em,itemsep=2pt]
\item À l'aide de $1+\tg^2\alpha=\dfrac{1}{\cos^2\alpha}$, détermine $\cos\alpha$ puis $\sin\alpha$
(coordonnées exactes de $M$).
\item Calcule $\sin2\alpha$ et $\cos2\alpha$.
\item Calcule $\tg2\alpha$ de deux façons (par $\dfrac{\sin2\alpha}{\cos2\alpha}$ et par la formule
$\tg2\alpha=\dfrac{2\tg\alpha}{1-\tg^2\alpha}$) et vérifie qu'elles coïncident.
\end{enumerate}

\exo{Une formule de l'angle triple}
On veut établir que, pour tout angle $a$, $\ \sin3a=3\sin a-4\sin^3a$.
\begin{enumerate}[label=\textbf{(\alph*)},leftmargin=2em,itemsep=2pt]
\item Écris $\sin3a=\sin(2a+a)$ et développe avec la formule d'addition.
\item Remplace $\sin2a$ et $\cos2a$ par leurs expressions, puis utilise $\cos^2a=1-\sin^2a$ pour
conclure.
\item Application : sachant $\sin a=\dfrac12$ (soit $a=30\degr$), vérifie que la formule redonne bien
$\sin90\degr=1$.
\end{enumerate}

\end{document}
